\section{Introduction}
\label{sec:introduction}

Product attribute/value pairs are crucial for e-commerce platforms since they enable online shoppers to use faceted product search~\cite{wei_survey_2013} and to compare products along explicit criteria. 
On e-commerce marketplaces and open product catalogues, product attribute/value pairs are often missing because merchants only provide unstructured product information, such as titles and descriptions~\cite{yang_mave_2022,zheng_opentag_2018}.
%Todo Alex (Done): Explain in which context they are noisy and have missing values? In product catalogs? Which and why? Maybe it is better to argue via large open marketplaces or plattforms on which such attributes are not avalable as merchants just submit textial product descriptions.
Product information extraction extracts attribute/value pairs from these product titles and product descriptions.
%Todo Alex (Done):  Same here: Do not use the terms "missing" and "recovering" but a proper definition of the term.
Existing works can be distinguished into closed extraction and open extraction approaches. Closed extraction assumes that the relevant attributes are known prior to the extraction ~\cite{yang_mave_2022, zheng_opentag_2018, zhu_multimodal_2020, xu_scaling_2019, wang_learning_2020}. Open extraction assumes that neither a complete attribute set nor all attribute values are known prior to the extraction~\cite{zhang_oa-mine_2022}. 
% Todo Alex (Done): Please use a definition that covers open and closed extraction as well, not only closed.
For example, consider the following title of a product offer \emph{``Canon EOS 1000D SLR 10 MP EF-S 18-55''}. A closed extraction approach extracts from this title an explicit attribute like \emph{``Resolution``} and returns the attribute/value pair \emph{``Resolution: 10 MP``}. An open extraction extracts all attribute/value pairs and returns a list of attribute/value pairs like \emph{``Brand: Canon``, ``Megapixels: 10``}.
%Can be removed as you also say this later. In this paper, we systematically evaluate how ChatGPT extracts attribute/value pairs from product titles. For this purpose, we experiment with closed and open prompts for product information extraction. 
%, because product titles are carefully designed by retailers~\cite{xu_scaling_2019} and consequently form a valuable source for missing attribute/value pairs.
%  Todo Alex (Done): "A user willing to extract" is not a realistic use case, refine.
%a user willing to extract a \emph{``Resolution``} attribute for the product. Closed product information extract aims at extracting the attribute/value pair \emph{``Resolution: 10 MP``} from the title.

% Todo Alex (Done): Do not use "we focus" as this is vague. Be concrete in what the paper does.

% Todo Alex (Done): Say "state of the art" models often rely on PMLs or provide a bit deeper overivew of history of product data extraction methods.
\begin{sloppypar}
Many state-of-the-art product information extraction methods~\cite{zhu_multimodal_2020, xu_scaling_2019, yang_mave_2022, zheng_opentag_2018}
%Todo Alex: Please cite some of these methods here!
rely on pre-trained language models (PLM) such as BERT.
The different approaches encode a product title using the PLM and add additional layers on top, which tag the sequence of an attribute value in the title.
% Todo Alex (Done): Also cover the open use case here. So no "requested"
OA-Mine is an example of an open extraction approach that relies on a PLM to generate attribute value candidates from product titles and to discover new attributes from the generated attribute value candidates~\cite{zhang_oa-mine_2022}.
The approaches have two main drawbacks: (1) they require training data to fine-tune the models, and (2) the fine-tuned models have problems generalizing to unseen attributes and attribute values.
% Todo Alex (Done): Also say that they have problems generalizing to unseen (or use term out-of-distribution) attributes (not only attribute values).
\end{sloppypar}

%Recent work has shown that scaling up the number of parameters and the training compute of language models can overcome drawbacks of models like BERT. 
Large autoregressive language models (LLMs) such as ChatGPT~\cite{ouyang_training_2022}, BLOOM~\cite{yong_bloom1_2023} or PaLM~\cite{chowdhery_palm_2022} have recently shown their potential to overcome these shortcomings. Due to being pre-trained on huge amounts of text as well as due to emergent effects resulting from the model size~\cite{wei_emergent_2022}, LLMs often demonstrate a better zero-shot performance compared to PLMs like BERT and are more robust to unseen examples~\cite{brown_language_2020}.

In this paper, we evaluate how ChatGPT extracts attribute/value pairs from product titles using prompt engineering and in-context learning
%Todo Alex (Done): Emergent abilities only refers to ability to do stuff that is not mentioned in training data. Here you need to also mention large training corpora.
and make the following contributions:
\begin{itemize}
    \item We evaluate the performance of ChatGPT (gpt3.5-turbo-0301) for closed and open product information extraction.
    \item We systematically explore different prompt designs and show the performance of the best zero-shot prompt design is similar to a fine-tuned pre-trained language model.
    \item We experiment with in-context learning for product information extraction and show that providing only three demonstrations is sufficient to reach an F1 score of 94.51\%.
    %\item We analyze the extracted attribute/value pairs and show how ChatGPT normalizes the attribute/value pairs.
    %Todo Alex: Remove the last contribution if you do not have space in the paper to propoerly present this.
\end{itemize}

This paper is structured as follows. Section \ref{sec:experimental_setup} introduces the dataset and the experimental setup. Section \ref{sec:prompt_design} and Section \ref{sec:in_context_learning} explain different prompt designs, in-context learning and present results. The code and the dataset to reproduce our results are online available\footnote{https://github.com/wbsg-uni-mannheim/pie\_chatgpt}.
